\documentclass[conference]{IEEEtran}

\usepackage{booktabs}
\usepackage{graphicx}
\usepackage{amsmath}
\usepackage{amssymb}
\usepackage{multirow}
\usepackage{xurl}
\usepackage{enumitem}
\usepackage{listings}
\usepackage{xcolor}
\usepackage{cite}

\lstset{
  basicstyle=\ttfamily\small,
  breaklines=true,
  frame=single,
  columns=fullflexible,
  keepspaces=true
}

\title{OracleGuard: Trust-Scored LLM-Based Test Oracle Generation\\with Differential Validation}

\author{
  \IEEEauthorblockN{Omkar Kabde}
  \IEEEauthorblockA{Chaitanya Bharathi Institute of Technology\\
  Hyderabad, India\\
  omkarkabde@gmail.com}
  \and
  \IEEEauthorblockN{Mohammed Imaduddin}
  \IEEEauthorblockA{Chaitanya Bharathi Institute of Technology\\
  Hyderabad, India\\
  mdimad005@gmail.com}
}

\begin{document}

\maketitle

\begin{abstract}
Any software test requires two items to be operational: inputs to work the program, and a criterion to determine whether the output is correct at all. Although automated input generation has grown by a long way with the methods of fuzzing, symbolic execution and search-based testing, the other half of this equation---building reliable test oracles---is a thorny stumbling block. To know whether output is right, one has to know what the program was meant to do, and that knowledge is usually stored in the head of the developer. This leads to slow, manual and more difficult to maintain creation of oracles as the code develops.

The issue has a number of interlocked dimensions. Formalizing informal expectations to make implementation assertions is part of the job, and even seasoned developers overlook edge cases and unspoken contracts. With time, change of code nullifies the already established assertions and transforms them into noise generators. And such systems, in most cases, never have their proper behavior documented at all.

Large Language Models provide a real chance in this case. They are able to read code, deduce intent based on signatures and docstrings, and generate syntactically plausible assertions. But plausible is not correct. Models hallucinate, miss boundary conditions, and give assertions that sound confident and are just plain wrong. Trusting LLM-generated oracles literally brings in the very type of silent mistakes that testing is expected to detect.

The idea behind OracleGuard is simple: generated assertions are supposed to be considered as candidates, not conclusions. In the framework, synthesis of assertions is preceded by LLM generation, and mutation-based differential testing and a multi-metric trust scoring system are utilized in a five-stage pipeline. Every candidate oracle earns its merit by proving its ability to detect faults that can be measured---making AI-aided oracle generation something that can be embraced with certainty.
\end{abstract}

\begin{IEEEkeywords}
test oracle generation, large language models, mutation testing, differential testing, AI-generated software engineering, automated testing, trust scoring
\end{IEEEkeywords}

% -------------------------------------------------
\section{Introduction}

\subsection{The Test Oracle Problem}

Software testing is based on two pillars: test inputs and test oracles. Automated test input generation has taken significant strides with the use of fuzzing, symbolic execution, and search-based methods~\cite{fraser2011evosuite,barr2015oracle}. The oracle problem, however, has turned out to be much more difficult to automate. It requires time and extensive domain knowledge to specify what a program should do, in detail, and on all the cases it covers---knowledge that developers rarely have. This has been one of the most crucial unsolved problems in software testing research~\cite{barr2015oracle,pezze2014automated}.

The challenge manifests itself in practice in a number of forms:

\begin{itemize}
\item \textbf{Manual work:} To write assertions that model real program behaviour, it is necessary to have both domain knowledge and sustained attention.
\item \textbf{Incompleteness:} Even well-written test suites lack edge cases and undocumented behavioral contracts.
\item \textbf{Maintenance burden:} Because code evolves, previously existing assertions tend to break in a manner difficult to tell whether this is due to a bug or not.
\item \textbf{Implicit specifications:} There is so much about what a program is supposed to do that is not written anywhere---it is present only in the mind of the developer.
\end{itemize}

\subsection{LLMs for Code Understanding}

Large Language Models have shown exceptionally good code comprehension. Present models can analyse program structure, derive intended behaviour based on signatures and documentation, and generate syntactically correct test assertions~\cite{jiang2024survey_codellm,davis2024understanding_llm_code,xu2022systematic}. This has inspired an ever-increasing literature on LLM-assisted testing~\cite{hayet2024chatassert,bodicoat2025understanding}, including assertion synthesis~\cite{hayet2024chatassert} and structured oracle generation~\cite{hossain2024togll}.

With that said, there are well-known failure modes of LLMs which are highly important for testing:

\begin{itemize}
\item \textbf{Hallucination:} A model can produce statements which appear right but test the wrong thing.
\item \textbf{Inconsistency:} The same code repeated can give different assertions, and there is no principled basis to favor either of them.
\item \textbf{No internal validation:} There is nothing about the architecture of a language model that asserts its outputs make sense.
\end{itemize}

Recent attempts have shown that these constraints can be avoided at least partially by fine-tuning. TOGLL~\cite{hossain2024togll} demonstrates that supervised ordered prompting with fine-tuning leads to better oracles and ones that are more effective bug detectors. Such work is devoted to the improvement of generation quality---the model becomes more effective at creating good oracles~\cite{hossain2024togll,hossain2025doc2oracll}. But high average quality leaves the most significant question open: when should a particular generated oracle actually be trusted? Good aggregate measures do not mean that a certain assertion can be relied upon. A validation mechanism executed at the level of particular oracles is a requisite for correctness-critical workflows~\cite{xiong2023express_confidence,liu2024uncertainty_estimation}.

\subsection{Our Contribution}

This gap is filled by OracleGuard. Instead of enhancing model training or prompt design, it uses the LLM as a black-box generator and proposes a validation-first architecture which evaluates generated assertions by mutation-based differential testing and multi-metric trust scoring before accepting them.

These are the major contributions:

\begin{enumerate}
\item \textbf{Hybrid architecture:} A validation layer wrapping LLM-based assertion generation in mutation-based differential testing.
\item \textbf{Multi-dimensional trust scoring:} A composite reliability score based on mutation kill rates, model confidence, consistency, coverage estimates, and complexity adjustments.
\item \textbf{Refinement engine:} Specific improvement recommendations based on surviving mutant patterns.
\item \textbf{Modular pipeline:} Five independently executable stages supporting research, debugging, and incremental development.
\item \textbf{Model-agnostic design:} Provider-level abstraction supporting a variety of LLM backends without task-specific fine-tuning.
\end{enumerate}

Where generation-centric methods such as TOGLL~\cite{hossain2024togll} are concerned with ensuring that the model generates better oracles, OracleGuard is concerned with choosing which oracles to trust. This framing---modeling outputs as requiring empirical acceptance---is consistent with new trends in LLM-assisted programming including verification, calibration of confidence, and trust-aware adoption~\cite{xiong2023express_confidence,liu2024uncertainty_estimation}.

% -------------------------------------------------
\section{Background and Related Work}

\subsection{Test Oracle Automation}

A test oracle is a device that decides how the behavior of a given program should be judged in response to an input~\cite{molina2025test,barr2015oracle,pezze2014automated}. The oracle problem---the difficulty of stating the correct behavior such that it can be checked automatically---has persisted as a central challenge~\cite{barr2015oracle}.

Automated methods of oracle construction are divided into a number of categories:

\begin{itemize}
\item \textbf{Assertion-based oracles:} executable predicates embedded directly in tests.
\item \textbf{Contract-based oracles:} pre/postconditions and invariants derived from formal specifications.
\item \textbf{Metamorphic oracles:} consistency relations across multiple executions~\cite{segura2016metamorphic,chen2018metamorphic}.
\item \textbf{Differential oracles:} cross-checking behavior between implementations or program variants.
\item \textbf{Trace-based oracles:} learned or rule-based classifiers over execution traces.
\end{itemize}

\subsection{LLM-Based Oracle and Test Generation}

Recent work proposes an increasing number of applications of LLMs to produce tests and oracles~\cite{hayet2024chatassert,bodicoat2025understanding,hossain2024togll}, including assertion synthesis, unit test generation, and translation of structured documentation to executable checks~\cite{hossain2024togll,hossain2025doc2oracll}. One thing that has been repeatedly observed is that the assertions produced by models can be useful but necessitate independent verification to prevent silent oracle errors~\cite{bodicoat2025understanding}. This fact is the direct basis of the design of OracleGuard.

\subsection{Mutation and Transformation-Based Validation}

Mutation testing compares the quality of tests by injecting small syntactic defects and quantifying the frequency of their detection by the tests~\cite{jia2011mutation,papadakis2019mutation}. It has become one of the standard measures of adequacy of test suites. The majority of previous literature utilizes mutation to assess available test suites~\cite{jia2011mutation,papadakis2019mutation}. OracleGuard does the opposite: mutation results are used to evaluate and score individual generated assertions prior to acceptance. The combination of mutation testing with LLMs has also been examined recently~\cite{straubinger2025mutation,zhang2025simulink,alimadadi2026hybrid}.

% -------------------------------------------------
\section{Motivation and Problem Definition}

An LLM is able to read a function, deduce what it is attempting to accomplish, and come up with assertions that appear plausible. But a statement may appear right and yet be silent in its inability to spot real bugs. It may fail to capture an important boundary case, or test the wrong thing entirely, providing developers with a false sense of security.

The ability to produce assertions does not inform us about which ones to believe, and this is the simple gap behind which OracleGuard is driven. The only way to bridge that gap is to consider generation and acceptance as quite distinct issues, each one supported by independent evidence.

\textbf{Problem statement.} Given a function under test with automatically generated inputs, synthesize executable assertions via an LLM, then generate a quantitative, empirically-based trust estimate depending on the fault-detection capability exhibited by the synthesized assertions---using a staged validation architecture.

\subsection{Validation Approach}

The validation method uses mutation-based differential testing~\cite{jia2011mutation,papadakis2019mutation} to stress-test candidate assertions. Various mutation operator classes are employed:

\begin{itemize}
\item Arithmetic operator substitutions
\item Relational operator changes
\item Logical operator transformations
\item Constant modifications
\item Statement deletions
\item Return-value mutations
\end{itemize}

All mutants are run against the generated test. When the test fails on the mutant, the mutation is killed; when the test passes, the mutation survives. Survival patterns indicate areas of weak assertions. The strength of the oracle can be measured empirically using aggregate kill rates~\cite{jia2011mutation}.

\begin{figure}[t]
\centering
\fbox{\parbox{0.92\linewidth}{
\textbf{Stage 1: Static Analysis} $\rightarrow$ metadata extraction, complexity analysis \\
\textbf{Stage 2: Test Prefix Generation} $\rightarrow$ fixture and input synthesis \\
\textbf{Stage 3: LLM Assertion Generation} $\rightarrow$ prompt construction and synthesis \\
\textbf{Stage 4: Differential Testing} $\rightarrow$ mutation generation and execution \\
\textbf{Stage 5: Analysis and Refinement} $\rightarrow$ trust scoring and recommendations
}}
\caption{OracleGuard five-stage pipeline architecture.}
\label{fig:og-pipeline}
\end{figure}

\subsection{Trust Quantification}

OracleGuard calculates a weighted composite trust score based on several sources of evidence~\cite{xiong2023express_confidence,liu2024uncertainty_estimation}:

\begin{itemize}
\item \textbf{Mutation score:} fraction of injected faults detected
\item \textbf{Model confidence:} average confidence across generated assertions
\item \textbf{Consistency score:} behavioral stability across mutant executions
\item \textbf{Coverage score:} estimated breadth of behavioral coverage
\item \textbf{Complexity penalty:} downward adjustment for high-complexity methods
\end{itemize}

The resultant score is scaled to the range $[0, 1]$ and translated to qualitative judgments---\texttt{VERIFIED}, \texttt{SUSPICIOUS}, \texttt{NEEDS\_REFINEMENT}---which provide a tangible foundation for acceptance, refinement, or rejection of a generated oracle.

\subsection{Research Questions}

\begin{itemize}
\item \textbf{RQ1:} Do mutation-based differential tests systematically expose weaknesses in LLM-generated assertions across a variety of mutation operator classes?
\item \textbf{RQ2:} Does a weighted multi-metric trust score offer a viable and understandable reliability indication for individual automatically generated oracles?
\end{itemize}

% -------------------------------------------------
\section{Approach}

\subsection{Architecture Overview}

OracleGuard performs a source-code-to-trusted-oracle pipeline in five stages. Stages communicate via typed intermediate artifacts, providing traceability and modular execution.

\subsection{Stage 1: Static Analysis}

The analysis stage derives method-level metadata from source code using Python's Abstract Syntax Tree (AST).

\textbf{Input:} Python source code.

\textbf{Process:}
\begin{enumerate}
\item Parse with \texttt{ast.parse}.
\item Walk the tree to determine function definitions.
\item Extract per-method metadata: signature and type annotations, parameter names and types, return type, docstring, dependency calls, and cyclomatic complexity.
\end{enumerate}

\textbf{Output:} Structured \texttt{MUTMetadata} objects. Methods are filtered optionally by complexity (e.g., 2--20 range).

\subsection{Stage 2: Test Prefix Generation}

This phase generates the setup code that is executed prior to the test method.

\textbf{Input:} \texttt{MUTMetadata}.

\textbf{Process:} Detect instance method or standalone function; create fixture instance when necessary; synthesize parameter values through one of three strategies: \textbf{Random} (type-appropriate random values), \textbf{Boundary} (edge-case values: 0, empty, limits), or \textbf{Equivalence} (representative partition values).

\begin{lstlisting}[language=Python]
int   -> random.randint(1, 100)
float -> random.uniform(0.0, 100.0)
str   -> f"test_{param_name}"
bool  -> random.choice([True, False])
list  -> [1, 2, 3] or []
\end{lstlisting}

\textbf{Output:} \texttt{TestPrefix} containing imports, setup code, and variable bindings.

\subsection{Stage 3: LLM Assertion Generation}

An LLM is used to synthesize candidate assertions, based on best practices in recent work~\cite{hayet2024chatassert,hossain2024togll,hossain2025doc2oracll}. A structured prompt with the function signature, full source code, docstring, generated inputs, invocation statement, and required JSON output schema is compiled and forwarded to the model.

\begin{lstlisting}
Generate assertions for this Python function.

Signature: def f(...)
Docstring: ...
Source: ...
Inputs: ...
Call: result = f(...)

Return JSON with assertions, explanation,
confidence, and type.
\end{lstlisting}

The provider layer is model-agnostic, supporting a number of LLM backends and a mock mode for offline testing. JSON responses are parsed to formatted assertions with fallback for format violations.

\textbf{Output:} \texttt{TestCase} with executable test code and candidate assertions.

\subsection{Stage 4: Differential Testing}

The generated assertions are stressed with mutation-based differential testing~\cite{jia2011mutation,papadakis2019mutation} using the following operator classes: arithmetic operator replacement, relational operator replacement, logical operator replacement, constant modification, statement deletion, and return-value mutation.

\textbf{Process:} Generate $N$ mutants via AST-level transformations; test each mutant in an independent subprocess with a time constraint; archive kill/survive data; compute the mutation score:

\[
\text{MutationScore} = \frac{|\text{Killed}|}{|\text{Killed}| + |\text{Survived}|}
\]

\textbf{Output:} \texttt{DifferentialReport} with per-mutant outcomes and aggregate statistics.

\subsection{Stage 5: Analysis and Refinement}

The final step combines signals into a trust-scored verdict, based on the principles of uncertainty quantification~\cite{xiong2023express_confidence,liu2024uncertainty_estimation}.

The weights in the trust metric are: mutation score (35\%), mean LLM confidence (20\%), consistency score (25\%), coverage score (15\%), and complexity penalty (5\%).

\[
T = 0.35M + 0.20L + 0.25C + 0.15V - 0.05P
\]

$T$ is brought to the range $[0, 1]$ and categorized as verdicts: \texttt{VERIFIED}, \texttt{SUSPICIOUS}, \texttt{NEEDS\_REFINEMENT}, or \texttt{REJECTED}. A refinement module takes the surviving mutants and generates specific recommendations.

\textbf{Output:} \texttt{OracleVerdict} that includes score, status, weaknesses, and recommendations.

% -------------------------------------------------
\section{Proposed Implementation Design}

\subsection{Design Principles and Extensibility}

The proposed implementation is structured on three principles: \textbf{Modularity} (each stage is independent with well-defined interfaces), \textbf{Extensibility} (interface-based components allow alternative models, operators, and input strategies), and \textbf{Reproducibility} (structured intermediate artifacts enable replay and systematic comparison).

The most important extension points are: \texttt{LLMProvider} (adapter interface to various backends), \texttt{MutationOperator} (pluggable transformation rules), and \texttt{InputStrategy} (interchangeable parameter generation strategies).

The prototype will be in Python, which has native AST support and subprocess isolation for safe mutant execution, as well as rich ecosystem support for LLM integration. External dependencies are introduced at the model integration level.

\subsection{Execution Model}

The pipeline is staged and sequential in nature, coordinated by a central orchestrator. Each stage produces a typed artifact which is the input of the next: method metadata, test prefix, candidate assertions, differential report, and trust-scored verdict. This enables partial execution, checkpoint replacement, and component substitution.

% -------------------------------------------------
\section{Evaluation}

\subsection{Why HumanEval+}

We needed a benchmark where the functions are self-contained, the correct implementations are known, and human-written test assertions exist for comparison. HumanEval+~\cite{liu2024evalplus}---an augmented version of OpenAI's HumanEval~\cite{chen2021humaneval} with 80$\times$ more test cases per problem---fits all three criteria. Its 164 Python problems require no external dependencies, which means our AST-based mutation engine can operate on them directly. And because the benchmark is widely used in LLM code generation research, results are comparable to prior work.

Other benchmarks were considered. Defects4J offers real bugs but targets Java. BugsInPy~\cite{widyasari2020bugsinpy} provides real Python bugs in production projects, but the functions have complex dependencies that our current prefix generator cannot handle---a natural extension for future work. MBPP has weaker ground-truth assertions (only three per problem), making it less useful as an oracle-quality reference.

We ran the evaluation on five HumanEval+ problems (IDs 0--4), covering list processing, string parsing, numeric computation, and sequence generation. Five problems is a small sample, and we do not claim statistical significance. What the sample does provide is a clear demonstration that the pipeline behaves consistently across function types and LLM providers.

\subsection{Setup}

Every model received the same prompt, the same temperature (0.7), and was evaluated through the same pipeline with 15 mutants per test case and one test case per problem. The only variable across runs was the model generating the assertions. Ten additional mutants, seeded independently of the ones used for scoring, were used to measure fault detection in a held-out fashion. Human-written HumanEval+ assertions served as the ground-truth baseline.

One detail matters for interpreting the results. We track two kinds of mutant kills. The \emph{total} kill rate counts any test failure---crashes included---as a kill. The \emph{oracle} kill rate counts only assertion failures; crashes from mutated code are caught by a try/except guard and treated as survived. The oracle kill rate isolates what the assertions themselves catch, and that is the signal ($M$) fed into the trust score.

\subsection{27 Models, 7 Providers}

We pointed OracleGuard at 27 different LLMs spanning OpenAI, OpenRouter, Groq, SambaNova, Cohere, NVIDIA, and Google Gemini. The models range from 1.2 billion parameters (Liquid LFM) to 671 billion (DeepSeek V3), and include both open-source and proprietary systems. Table~\ref{tab:models} shows the results.

\begin{table}[t]
\centering
\caption{Model ranking by mean trust score (top 10 and bottom 3 of 27). $^*$Reasoning model.}
\label{tab:models}
\small
\begin{tabular}{rlccccccc}
\toprule
\# & \textbf{Model} & \textbf{Provider} & \textbf{Trust} & \textbf{V} & \textbf{S} & \textbf{N} & \textbf{R} \\
\midrule
1 & Mistral 7B v0.2 & NVIDIA & 0.598 & 2 & 1 & 0 & 2 \\
2 & Gemini 3.1 FL & Gemini & 0.552 & 1 & 2 & 0 & 2 \\
3 & GPT-4.1 Mini & OpenAI & 0.547 & 2 & 0 & 1 & 2 \\
4 & GPT-OSS 20B & OpenRouter & 0.545 & 1 & 1 & 1 & 2 \\
5 & GPT-4o & OpenAI & 0.500 & 2 & 0 & 0 & 3 \\
6 & DeepSeek V3 & SambaNova & 0.486 & 1 & 1 & 0 & 3 \\
7 & Kimi K2 & Groq & 0.482 & 1 & 1 & 0 & 3 \\
8 & Llama 3.1 8B & Groq & 0.478 & 2 & 0 & 0 & 3 \\
9 & Llama 3.1 8B & SambaNova & 0.478 & 2 & 0 & 0 & 3 \\
10 & Devstral 123B & NVIDIA & 0.478 & 1 & 1 & 0 & 3 \\
\midrule
25 & Gemini 3.0 FL$^*$ & Gemini & 0.273 & 0 & 1 & 0 & 4 \\
26 & Qwen3 32B$^*$ & Groq & 0.136 & 0 & 1 & 0 & 4 \\
27 & DeepSeek R1$^*$ & SambaNova & 0.000 & 0 & 0 & 0 & 5 \\
\bottomrule
\end{tabular}
\end{table}

A few things stand out immediately. Mistral 7B---a seven-billion-parameter open-source model hosted on NVIDIA---leads the ranking. GPT-4o, a far larger proprietary model, sits at fifth place. And the three reasoning models (DeepSeek R1, Qwen3 32B, Gemini 3.0 Flash) cluster at the bottom, not because they write worse assertions, but because their \texttt{<think>} tags break JSON extraction. We return to this in the discussion.

\subsection{Does the Trust Score Actually Predict Fault Detection?}

This is the question that matters. A trust score that correlates with nothing is decoration.

We ran a controlled experiment separate from the benchmark. Four sets of synthetic assertions---strong (exact value checks), medium (type plus range), weak (not-null only), and bad (tautologies like \texttt{assert True})---were applied to five functions with known behavior. For each set, we computed the trust score and then independently measured fault detection on fresh mutants.

The trust score tracks fault detection closely. Spearman's rank correlation is $\rho = 0.738$---a strong positive relationship. The ordering is monotonic: strong oracles average 0.334 trust and 10.4\% fault detection; bad oracles average 0.127 trust and 0\% fault detection. The gap between the two (0.207) is large enough to be useful in practice.

\subsection{OracleGuard vs.\ Ground Truth}

Across all 27 models, OracleGuard's assertions achieve 100\% fault detection on seeded mutants when crash-kills are included. Human-written ground-truth assertions from HumanEval+ catch 84\%. That gap is partly an artifact---mutations often crash the function before any assertion runs, and OracleGuard's test harness counts crashes as kills. But even after adjusting for this, the generated assertions perform respectably against a hand-crafted baseline.

The per-operator breakdown tells a more nuanced story. Return-value mutations and arithmetic swaps are caught by virtually every model (100\% kill rate). Constant modifications (65--100\%) and relational operator changes (50--100\%) show the widest variation. The pattern is predictable: models that write \texttt{assert result == 80.0} catch constant perturbations; models that write only \texttt{assert isinstance(result, float)} do not.

\subsection{Provider Independence}

We ran Llama 3.1 8B on both Groq and SambaNova. Same model weights, different inference infrastructure. The trust score came back 0.478 on both, with identical verdict distributions (2V/0S/0N/3R). OracleGuard evaluates the oracle, not the provider.

% -------------------------------------------------
\section{Discussion}

\subsection{Size Does Not Predict Quality}

If there is one result we did not expect, it is this: the number of parameters in a model tells you almost nothing about how good its test assertions will be.

\begin{itemize}
\item Mistral 7B (trust 0.598) outperforms DeepSeek V3 at 671 billion parameters (0.486) and Devstral at 123B (0.478).
\item Liquid LFM, a 1.2B model, scores 0.446---ahead of Nemotron at 30B (0.380) and Qwen3 at 32B (0.136).
\item Command A at 111B (0.400) scores lower than Command R at 35B (0.453).
\end{itemize}

What seems to matter instead is how precisely the model follows the output schema. Models that reliably produce well-formed JSON with exact expected values (``\texttt{assert result == 80.0}'') score well. Models that pad their output with verbose explanations or fall back to generic type checks (``\texttt{assert isinstance(result, float)}'') do not. For teams choosing an LLM for oracle generation, this means smaller and cheaper models can work just as well---possibly better---than the largest available option.

\subsection{Reasoning Models and the Output Format Problem}

The bottom three slots in the ranking belong to reasoning models: DeepSeek R1 (0.000), Qwen3 32B (0.136), and Gemini 3.0 Flash (0.273). All three wrap their responses in \texttt{<think>...</think>} tags that precede the actual JSON. Our parser strips these tags, but in many cases the JSON itself is truncated or malformed by the time the thinking is done.

This is not a reflection of assertion quality. It is an output format problem. The one reasoning model that performs adequately---o4-mini at 0.465---does so because OpenAI strips thinking tokens server-side before returning the API response. The lesson is practical: reasoning models need output post-processing tailored to their format, not just better prompts.

\subsection{Run-to-Run Variance}

We ran Kimi K2 five times on the same five problems. The trust scores were: 0.443, 0.404, 0.360, 0.350, 0.482. That is a 37\% coefficient of variation from a model that did not change between runs.

The explanation is straightforward. At temperature 0.7, the model sometimes generates a strong value assertion and gets VERIFIED. Other times it generates only a type check and gets REJECTED. Both scores are correct---they reflect the quality of what was actually generated in each run. But they also mean that a single evaluation pass is not enough to characterize a model. We recommend running at least three times and reporting the mean and standard deviation.

\subsection{What Surviving Mutants Tell You}

The trust score gives you a number. The surviving mutant log gives you something you can act on.

A score of 0.69 (SUSPICIOUS) is useful for policy---should this oracle go into the test suite? But ``five constant-replacement mutants survived, specifically mutations changing \texttt{100} to \texttt{101} on line~7'' is useful for \emph{fixing} the oracle. The developer knows the assertion set lacks an exact-value check that would break under constant perturbation.

The mapping from surviving operators to missing assertions is consistent across models:

\begin{itemize}
\item Relational survivors $\rightarrow$ add boundary comparison assertions
\item Constant survivors $\rightarrow$ tighten value equality checks
\item Statement-deletion survivors $\rightarrow$ add coverage for removed code paths
\item Return-value survivors $\rightarrow$ add explicit type-and-value checks
\end{itemize}

These suggestions can be fed back to the LLM as a refinement prompt: ``The following mutants survived your assertions. Generate additional assertions that would detect these specific faults.'' In early experiments, this loop converges in two to three rounds for most functions.

\subsection{Complement to Generation-Focused Work}

TOGLL~\cite{hossain2024togll} and similar systems ask: how do we make the model generate better oracles on average? OracleGuard asks a different question: for \emph{this particular} oracle, should we trust it? Both questions matter. A team that runs TOGLL's fine-tuned generator through OracleGuard's validation layer gets the best of both---higher average quality from better generation, and per-oracle trust signals from mutation-based validation. The two lines of work are not in competition. They are two halves of the same pipeline.

% -------------------------------------------------
\section{Threats to Validity}

\paragraph{Internal.} Five problems evaluated once per model is enough to show the pipeline works, but not enough to make strong statistical claims. LLM non-determinism is a real confound---Kimi K2 varied 37\% across five runs on the same inputs. The trust score weights (0.35, 0.20, 0.25, 0.15, 0.05) were validated through correlation analysis ($\rho = 0.738$) but not calibrated on a held-out dataset. Different weight choices could shift borderline verdicts.

\paragraph{External.} We tested on self-contained Python functions from HumanEval+. Production code has dependencies, state, file I/O, and concurrency---none of which our prefix generator or mutation engine currently handles. The random input generator sometimes produces type-mismatched values (a single float where a list is expected), which depresses trust scores on certain problems. BugsInPy~\cite{widyasari2020bugsinpy} would provide a more realistic evaluation target.

\paragraph{Construct.} Mutation score is a well-established proxy for test quality~\cite{jia2011mutation,papadakis2019mutation}, but it is not the same thing as correctness. Our mutations apply to the entire source file rather than the target function alone, so some mutants are irrelevant to the test being evaluated. We do not handle equivalent mutants, which inflate the denominator and depress scores. The oracle kill rate---counting only assertion failures, not crashes---mitigates this partially but does not eliminate it.

% -------------------------------------------------
\section{Conclusion and Future Work}

OracleGuard started from a simple premise: treat every LLM-generated assertion as a hypothesis, not a conclusion, and let mutation testing decide which hypotheses hold up. The five-stage pipeline that implements this premise---static analysis, prefix generation, LLM synthesis, differential testing, and trust-scored refinement---turns oracle generation from a one-shot gamble into something measurable.

We evaluated the framework across 27 models from seven providers. The trust score correlates strongly with actual fault detection ($\rho = 0.738$), and it separates good oracles from bad ones with a gap wide enough to be useful in practice (0.207 between strong and tautological assertions). Model size turned out to be a poor predictor of oracle quality; a 7B model led the ranking over systems 100$\times$ larger. What mattered was whether the model could follow the output schema and produce precise expected values.

The most practically useful output is not the score---it is the surviving mutant log. A number like 0.69 informs a policy decision. A report saying ``three constant-replacement mutants survived on line~7'' informs a developer decision. That specificity is what turns validation into refinement.

\subsection{What Comes Next}

Several things would strengthen the work and broaden its reach:

\begin{itemize}
\item \textbf{Scale.} Running on 50+ HumanEval+ problems with multiple passes per model would move the results from directional to statistically significant.
\item \textbf{Real bugs.} BugsInPy~\cite{widyasari2020bugsinpy} offers 493 real bugs in production Python projects. Validating against real faults, not just seeded mutations, would test whether the trust score generalizes beyond synthetic defects.
\item \textbf{Weight calibration.} The trust score weights are design choices. A held-out validation set and sensitivity analysis would either confirm them or suggest better alternatives.
\item \textbf{Closing the loop.} Right now Stage~5 suggests improvements and stops. Feeding those suggestions back to Stage~3---prompting the LLM with the surviving mutant report---would turn the pipeline into an iterative system. Early experiments suggest convergence in two to three rounds.
\item \textbf{Reasoning models.} Stripping \texttt{<think>} tags is a partial fix. Robust structured-output extraction from reasoning models is an open problem for any system consuming LLM output, not just ours.
\end{itemize}

OracleGuard does not write perfect oracles. No system will. But it tells you, clearly and specifically, which oracles are weak and why---and that is something developers have not had before.

% -------------------------------------------------
\bibliographystyle{IEEEtran}
\bibliography{references}

% -------------------------------------------------
\appendix

\section{Worked Example}

\subsection{Source Function}

\begin{lstlisting}[language=Python]
def calculate_discount(price: float,
                       discount_percent: float) -> float:
    """Calculate discounted price. Raises ValueError
    if discount_percent is outside [0, 100]."""
    if discount_percent < 0 or discount_percent > 100:
        raise ValueError(
            "Discount must be between 0 and 100")
    discount_amount = price * (discount_percent / 100)
    return round(price - discount_amount, 2)
\end{lstlisting}

\subsection{Example Generated Test}

\begin{lstlisting}[language=Python]
def test_calculate_discount_0():
    arg_price = 100.0
    arg_discount_percent = 20.0
    result = calculate_discount(
        arg_price, arg_discount_percent)
    assert result is not None
    assert isinstance(result, float)
    assert result == 80.0
    assert 0 <= result <= arg_price
\end{lstlisting}

\subsection{Example Mutation Validation Outcome}

Under mutation-based validation, arithmetic and value mutations are typically killed by these assertions---the \texttt{result == 80.0} check catches most arithmetic operator substitutions and constant modifications. However, mutations removing the \texttt{ValueError} guard for out-of-range inputs survive, since no assertion checks exception behavior. The framework surfaces this as a surviving mutant in the \texttt{STATEMENT\_DELETE} class and recommends adding an exception assertion for boundary inputs. This outcome feeds directly into the trust score and refinement suggestions presented to the developer.

\section{Full Model Ranking}

Table~\ref{tab:full-ranking} presents the complete ranking of all 27 models evaluated. Models are sorted by mean trust score descending. Size indicates publicly known parameter count where available. Reasoning models are marked with $^*$.

\begin{table*}[t]
\centering
\caption{Complete model ranking by mean trust score across 5 HumanEval+ problems.\\V = VERIFIED, S = SUSPICIOUS, N = NEEDS\_REFINEMENT, R = REJECTED.}
\label{tab:full-ranking}
\small
\begin{tabular}{rllrccccccc}
\toprule
\# & \textbf{Model} & \textbf{Provider} & \textbf{Size} & \textbf{OSS} & \textbf{Trust} & \textbf{V} & \textbf{S} & \textbf{N} & \textbf{R} \\
\midrule
1 & Mistral 7B v0.2 & NVIDIA & 7B & Y & 0.598 & 2 & 1 & 0 & 2 \\
2 & Gemini 3.1 Flash Lite & Gemini & --- & N & 0.552 & 1 & 2 & 0 & 2 \\
3 & GPT-4.1 Mini & OpenAI & --- & N & 0.547 & 2 & 0 & 1 & 2 \\
4 & GPT-OSS 20B & OpenRouter & 20B & Y & 0.545 & 1 & 1 & 1 & 2 \\
5 & GPT-4o & OpenAI & --- & N & 0.500 & 2 & 0 & 0 & 3 \\
6 & DeepSeek V3 & SambaNova & 671B & Y & 0.486 & 1 & 1 & 0 & 3 \\
7 & Kimi K2 & Groq & --- & Y & 0.482 & 1 & 1 & 0 & 3 \\
8 & Gemini 2.5 Flash Lite & Gemini & --- & N & 0.478 & 2 & 0 & 0 & 3 \\
9 & Llama 3.1 8B & Groq & 8B & Y & 0.478 & 2 & 0 & 0 & 3 \\
10 & Llama 3.1 8B & SambaNova & 8B & Y & 0.478 & 2 & 0 & 0 & 3 \\
11 & Devstral 123B & NVIDIA & 123B & Y & 0.478 & 1 & 1 & 0 & 3 \\
12 & o4-mini$^*$ & OpenAI & --- & N & 0.465 & 1 & 1 & 0 & 3 \\
13 & Qwen 2.5 Coder 7B & NVIDIA & 7B & Y & 0.454 & 1 & 1 & 0 & 3 \\
14 & Command R & Cohere & 35B & N & 0.453 & 1 & 1 & 0 & 3 \\
15 & Liquid LFM 1.2B & OpenRouter & 1.2B & Y & 0.446 & 2 & 0 & 0 & 3 \\
16 & GPT-5.1 & OpenAI & --- & N & 0.441 & 0 & 2 & 0 & 3 \\
17 & Falcon3 7B & NVIDIA & 7B & Y & 0.440 & 1 & 1 & 0 & 3 \\
18 & GPT-5.4 & OpenAI & --- & N & 0.422 & 1 & 0 & 1 & 3 \\
19 & Minimax M2.5 & OpenRouter & --- & N & 0.408 & 1 & 0 & 1 & 3 \\
20 & Command A & Cohere & 111B & N & 0.400 & 0 & 1 & 1 & 3 \\
21 & Gemma 3 4B & OpenRouter & 4B & Y & 0.393 & 1 & 0 & 1 & 3 \\
22 & Nemotron 30B & OpenRouter & 30B & N & 0.380 & 2 & 0 & 0 & 3 \\
23 & Qwen 3.6+ & OpenRouter & --- & N & 0.314 & 0 & 0 & 2 & 3 \\
24 & Phi-3 Medium & NVIDIA & 14B & Y & 0.308 & 0 & 1 & 0 & 4 \\
25 & Gemini 3.0 Flash$^*$ & Gemini & --- & N & 0.273 & 0 & 1 & 0 & 4 \\
26 & Qwen3 32B$^*$ & Groq & 32B & Y & 0.136 & 0 & 1 & 0 & 4 \\
27 & DeepSeek R1$^*$ & SambaNova & 671B & Y & 0.000 & 0 & 0 & 0 & 5 \\
\bottomrule
\end{tabular}
\end{table*}

\end{document}
